% !TEX root = ../matan.tex

\section{Теорема о почленном дифференцировании функционального ряда.}

\begin{Theorem}{Почленное дифференцирование}{}
	Пусть $I = [a, b]$, функции $f_k \colon I \to \RR$ непрерывно дифференцируемы на $I$, и
	\begin{enumerate}
		\item в некоторой точке $x_0 \in I$ числовой ряд $\sum_{k=1}^{\infty} f_k(x_0)$ сходится;
		\item ряд производных $\sum_{k=1}^{\infty} f_k'$ сходится равномерно на $I$ к функции $g$.
	\end{enumerate}
	Тогда ряд $\sum_{k=1}^{\infty} f_k$ сходится равномерно на $I$ к некоторой функции $f$, которая дифференцируема, причём
	\[
	f'(x) = g(x) = \sum_{k=1}^{\infty} f_k'(x), \qquad \text{то есть} \quad \Bigl( \sum_{k=1}^{\infty} f_k(x) \Bigr)' = \sum_{k=1}^{\infty} f_k'(x).
	\]
\end{Theorem}

\begin{proof}
	Обозначим $S_n(x) = \sum_{k=1}^{n} f_k(x)$ и $G_n(x) = \sum_{k=1}^{n} f_k'(x)$; по условию $G_n \to g$ равномерно на $I$.

	\textbf{Равномерная сходимость $\sum f_k$.} Так как $S_n$ непрерывно дифференцируема и $S_n' = G_n$, по формуле Ньютона--Лейбница для $m > n$ и $x \in I$
	\[
	S_m(x) - S_n(x) = \bigl(S_m(x_0) - S_n(x_0)\bigr) + \int_{x_0}^{x} \bigl(G_m(t) - G_n(t)\bigr)\,dt,
	\]
	откуда
	\[
	|S_m(x) - S_n(x)| \le |S_m(x_0) - S_n(x_0)| + (b - a)\,\|G_m - G_n\|_\infty.
	\]
	Первое слагаемое мало, поскольку числовой ряд $\sum f_k(x_0)$ сходится (его частичные суммы фундаментальны); второе мало, поскольку $\{G_n\}$ равномерно фундаментальна (ряд $\sum f_k'$ равномерно сходится). Значит, $\sup_{x \in I} |S_m(x) - S_n(x)| \to 0$ при $m, n \to \infty$, и по критерию Коши $S_n \to f$ равномерно на $I$.

	\textbf{Дифференцируемость и вид $f$.} По формуле Ньютона--Лейбница
	\[
	S_n(x) = S_n(x_0) + \int_{x_0}^{x} G_n(t)\,dt.
	\]
	Переходим к пределу при $n \to \infty$: слева $S_n(x) \to f(x)$ и $S_n(x_0) \to \sum_k f_k(x_0)$, а переход под знак интеграла законен ввиду равномерной сходимости $G_n \to g$:
	\[
	\Bigl| \int_{x_0}^{x} (G_n(t) - g(t))\,dt \Bigr| \le (b - a)\,\|G_n - g\|_\infty \to 0.
	\]
	Получаем
	\[
	f(x) = \sum_{k=1}^{\infty} f_k(x_0) + \int_{x_0}^{x} g(t)\,dt.
	\]
	Функция $g$ непрерывна как равномерный предел непрерывных $G_n$, поэтому по основной теореме анализа $x \mapsto \int_{x_0}^{x} g(t)\,dt$ дифференцируема с производной $g(x)$. Следовательно, $f$ дифференцируема и
	\[
	f'(x) = g(x) = \sum_{k=1}^{\infty} f_k'(x). \qedhere
	\]
\end{proof}

\begin{Remark}{Зачем нужна сходимость в одной точке}{}
	Одной лишь равномерной сходимости ряда производных недостаточно: она определяет $f$ лишь с точностью до аддитивной постоянной (производные не различают функции, отличающиеся на константу). Условие сходимости $\sum f_k(x_0)$ в одной точке фиксирует эту постоянную и гарантирует сходимость самого ряда $\sum f_k$.
\end{Remark}
